\documentclass[aspectratio=169, 12pt]{beamer}
\usepackage{graphicx}
\usepackage{hyperref}
\hypersetup{
    colorlinks=true,
    urlcolor=MidnightBlue,
    linkcolor=MidnightBlue,
    citecolor=MidnightBlue
}
\usepackage{booktabs}
\usepackage{enumerate}
\usepackage[absolute,overlay]{textpos}
\usepackage{amsfonts}
\usepackage{colortbl}
\usepackage{tcolorbox}
\usepackage[default]{lato}
\usepackage[dvipsnames]{xcolor}
\usepackage{pgfplots}
\pgfplotsset{compat=1.16}

\setbeamertemplate{itemize items}[default]
\setbeamertemplate{section in toc}[circle]
\setbeamertemplate{subsection in toc}[circle]
\setbeamertemplate{navigation symbols}{}
\setbeamertemplate{caption}[numbered]
\usepackage{comment}
\usepackage{textcomp}

\newcommand{\backupbegin}{
   \newcounter{finalframe}
   \setcounter{finalframe}{\value{framenumber}}
}
\newcommand{\backupend}{
   \setcounter{framenumber}{\value{finalframe}}
}

\AtBeginSection[]{
  \begin{frame}
  \vfill
  \centering
    \usebeamerfont{title}\usebeamercolor[fg]{title} \LARGE \insertsectionhead\par%
  \vfill
  \end{frame}
}

\newcommand{\hlcode}[1]{\tcbox[colback=orange!10, colframe=orange!10, arc=3pt, boxrule=0pt, left=2pt, right=2pt, top=1pt, bottom=1pt, on line]{\texttt{#1}}}

\setbeamerfont{frametitle}{size=\large}
\setbeamertemplate{footline}[frame number]
\setbeamersize{text margin left=4mm, text margin right=4mm}

% ---------- URL shortcuts ----------
\newcommand{\ghLtwo}{https://tannerregan.github.io/GWU_AI_study_group/code/lecture_2/}
% -----------------------------------

\begin{document}
\linespread{1.3}

\title{Session 2: AI for Coding --- \\Practical Workflows and Demos}
\author[AI for Economists]{Elira Kuka and Tanner Regan}
\date[\today]{The George Washington University}
\maketitle


%%%%%%%%%%%%%%
\begin{frame}
\frametitle{Plan for Today: hands-on workflows}

\begin{enumerate}
    \item Demos with Stata:
    \begin{itemize}
        \item[\circ] Replication package {\color{gray}\it EK}
        \item[\circ] Checking RA work {\color{gray}\it EK}
        \item[\circ] Quick demo of with VS Code {\color{gray}\it TR}
    \end{itemize}
    \vspace{0.05in}
    \item Python and Claude Code
    	\begin{itemize}
    		\item[\circ] Python in VS Code overview  {\color{gray}\it TR}
    		\begin{itemize}
        			\item libraries and environments
        			\item Jupyter extension 
    		\end{itemize}
    		\item[\circ] Data analysis demo with Python
    		\begin{itemize}
        			\item Python example with Claude Code from VS Code
    		\end{itemize}
	\end{itemize}
\end{enumerate}

\end{frame}


\section{\textcolor{violet}{Coding with Claude Code \vspace{0.3in}}}


%%%%%%%%%%%%%%
\begin{frame}
\frametitle{Demo: Replication Package}

\begin{itemize}
    \item Task: take a paper written in Latex and code folder, and produce repl package
    \item Steps we will walk through:
        \begin{enumerate}
            \item Point Claude at the project folder
            \item Tell Claude to follow AEA instructions for replication package
            \item Tell Claude to test all code runs and produces correct output
            \item Save steps for future work (save as skill)
        \end{enumerate}
    \item What to watch for:
        \begin{itemize}
            \item[$\circ$] Bad prompting (mine!)
             \item[$\circ$] More inefficient that I hoped, though MUCH faster than myself
        \end{itemize}
\end{itemize}

\end{frame}


%%%%%%%%%%%%%%
\begin{frame}
\frametitle{Demo: Check RA work}

\begin{itemize}
    \item Task: undertand RA work and check for issues
        \begin{enumerate}
            \item Ask Claude to write a plan before touching any files
            \item Review the plan in Markdown; request changes to the approach
            \item Approve; Claude executes and produces a clean, commented script
            \item Compare output to Ask-mode version: same result, better code
        \end{enumerate}
    \item Reviewing session history: how to audit what Claude did
    \item Ask Claude to update README with what was done (good practice)
\end{itemize}

\end{frame}


%%%%%%%%%%%%%%
\begin{frame}
\frametitle{How to Verify AI Code}

\begin{itemize}
    \item AI code can look completely correct and be wrong in subtle ways
    \item For replication package, easy to check: results either replicates or not.
    \item Verification checklist for data work;
        \begin{itemize}
            \item[$\circ$] Does the code run without errors? (necessary, not sufficient)
            \item[$\circ$] Do merged/unmerged counts make sense at each merge step? 
       		 \begin{itemize}
           		\item[$\circ$] I ask my RA to do this
			\item[$\circ$] I am sure can ask Claude to provide these numbers
        		\end{itemize}
            \item[$\circ$] Do obdservation counts and means match codebook or prior literature?
            \item[$\circ$] Can you reproduce a known result (e.g., a published Table 1)?
            \item[$\circ$] Ask Claude: ``what assumptions did you make?'' --- it will often reveal them
        \end{itemize}
\end{itemize}

\end{frame}



%%%%%%%%%%%%%%
\begin{frame}
\frametitle{Stata in VS Code demo}

\begin{itemize}
    \item Very quickly - how to run stata from VS Code
    \begin{itemize}
        \item Use the \href{https://github.com/tmonk/stata-workbench}{Stata Workbench Extension} (setup guide at link)
        \item Make sure the extension can `find' stata
        \item Work (mostly) as normal
        \begin{itemize}
            \item Some things different, e.g. 'browse' needs to be done point and click.
        \end{itemize}
    \item Otherwise, Claude Code can work with Stata similarly to previous example
    \end{itemize}
    \textcolor{gray}{\textit{$\triangleright$ Stata \& VS Code Demo}}
\end{itemize}

\end{frame}



\section{\textcolor{violet}{Python in VS Code overview \vspace{0.3in}}}

%%%%%%%%%%%%%%
\begin{frame}
\frametitle{Why Python?}


\begin{itemize}
    \item \href{https://claesbackman.com/claude-code-guide.html\#:~:text=of\%20Economics.-,Python,-ms\%2Dpython.python}{Claes Backman}
\end{itemize}
Python for Economists: Common Tasks

\begin{itemize}
    \item Data analysis with \texttt{pandas}: merging, collapsing, reshaping --- analogous to Stata
    \item Econometrics: \texttt{statsmodels}, \texttt{linearmodels} (panel data, IV)
    \item Data collection: \texttt{requests}, \texttt{BeautifulSoup} for scraping (Session 3)
    \item Figures: \texttt{matplotlib}, \texttt{seaborn} --- publication-quality plots
    \vspace{.05in}
    \item Claude Code can translate your Stata code to Python (or vice versa) --- useful for:
        \begin{itemize}
            \item[$\circ$] Running analyses you cannot do in Stata
            \item[$\circ$] Working with large datasets where Python is faster
        \end{itemize}
\end{itemize}
\end{frame}


%%%%%%%%%%%%%%
\begin{frame}
\frametitle{Python basics}
\begin{itemize}
    \item Libraries and Environments
    \begin{itemize}
        \item Anaconda manages Python installations and packages
        \item Each project gets its own \textbf{environment}: a fixed set of packages and versions
            \begin{itemize}
                \item[$\circ$] Prevents conflicts between projects
                \item[$\circ$] Ensures others can reproduce your results exactly
            \end{itemize}
        \item Ask Claude Code to create and activate an environment for you
    \end{itemize}

    \item Three ways to run Python code
        \begin{itemize}
            \item[$\circ$] \textbf{Script} (\texttt{.py}): run the full pipeline --- best for production code
            \item[$\circ$] \textbf{Interactive window}: run \texttt{.py} scripts cell-by-cell --- best of both worlds
            \item[$\circ$] \textbf{Jupyter notebook} (\texttt{.ipynb}): interactive cells, inline output --- good for exploration
        \end{itemize}
\end{itemize}
\end{frame}




%%%%%%%%%%%%%%
\begin{frame}
\frametitle{Python Setup in VS Code}
\begin{itemize}
    \item Required extensions: Python, Jupyter
    \item Managing environments:
        \begin{itemize}
            \item[$\circ$] Each project gets its own environment (set of packages + versions)
            \item[$\circ$] Prevents conflicts between projects
            \item[$\circ$] Ask Claude Code to create and activate an environment for you
        \end{itemize}
    \item Running code:
        \begin{itemize}
            \item[$\circ$] \textbf{Jupyter notebook} (\texttt{.ipynb}): interactive cells, inline output --- good for exploration
            \item[$\circ$] \textbf{Jupyter interactive window}: run \texttt{.py} scripts cell-by-cell --- best of both worlds
            \item[$\circ$] \textbf{Script} (\texttt{.py}): run the full pipeline --- best for production code
        \end{itemize}
\end{itemize}
\end{frame}


\section{Data Analysis Demo with Python \vspace{0.3in}}


%%%%%%%%%%%%%%
\begin{frame}
    \begin{center}
        \color{red}\bf I AM IMPRESSED
    \end{center}
    \pause
    \begin{itemize}
        \item This example took about 2.5 days of my work
        \item Before Claude Code:  
        \begin{itemize}
            \item 2-4 weeks of my time
            \item 1-2 months of very good PhD RA time
            \item 2+ months of decent MS RA time
        \end{itemize}
        \item Claude solution still more efficient and elegant
        \item More discussion later...
    \end{itemize}
\end{frame}


%%%%%%%%%%%%%%
\begin{frame}
\frametitle{Python Exercise}

Produced data (sampling frame for survey) and a \href{\ghLtwo/python_demo/4_slides_event_study.pdf}{set of slides} in four steps:
\begin{enumerate}
    \item Download data from Kigali City
    \item Download data from Alpha Earth
    \item Build a sample frame with visualization
    \item Test an event study design with google satellite embeddings
\end{enumerate}
\vfill\pause
\begin{itemize}
    \item Paul GP has 3 complementary blog posts w/ youtube lectures
    \begin{itemize}
        \item \href{https://paulgp.substack.com/p/from-an-empty-folder-to-a-figure}{Data Analysis}, \href{https://paulgp.substack.com/p/from-edgar-filings-to-a-structured}{Web Scraping}, \href{https://paulgp.substack.com/p/large-datasets-and-structured-databases}{Large Datasets}
    \end{itemize}
\end{itemize}
\end{frame}

%%%%%%%%%%%%%%
\begin{frame}
\frametitle{Download data from Kigali City}

\begin{enumerate}
    \item Wrote \href{\ghLtwo/python_demo/Prompt data scraping for Kigali}{Prompt data scraping for Kigali.md} pasted into ``plan mode''
    \begin{itemize}
        \item I had to give permission on a few steps, e.g. searching URLs.
    \end{itemize}
    \item Claude generated the plan \href{\ghLtwo/python_demo/kigali-city-data-mutable-frost}{kigali-city-data-mutable-frost.md} ($\sim$10 mins)
    \item Executed the plan which created the script \href{\ghLtwo/python_demo/download_kigali_gis_v0.py}{download\_kigali\_gis\_v0.py}

    \item Ran the code without edits (using Jupyter window). Two issues: 
    \begin{itemize}
        \item Claude added packages to the yml but didn't update the conda environment. Chatted with Claude to identify and fix this. %Told claude to update the project CLAUDE.md file so that it always: (1) updates the python environment with newly required packages, and (2) updates the .yml file so that the environment is perfectly replicable.
        \item I had daycare pickup and the code was still running. Claude adjusted the code so it would restart where I left off, rather than restarting completely.
%- I told Claude "the script is running. I need to close my laptop and go home. Can I pause the script and restart it without disrupting the download?"
%-CC: "The script has no resume logic — if you stop and restart it, it will re-download everything from scratch, including layers already completed. The fix is simple: add a file-exists check so it skips layers that are already on disk. Want me to add that now before you leave? It's a two-line change per download function and would make the restart seamless."
    \end{itemize}
    \item Ran the updated script \href{\ghLtwo/python_demo/download\_kigali\_gis.py}{download\_kigali\_gis.py} at home $\rightarrow$ success!
    \begin{itemize}
        \item VS Code diff: right-click script.py ``Select for Compare'', right-click script2.py ``Compare with Selected''. Gives a clean side-by-side diff in the editor.
    \end{itemize}
\end{enumerate}
\end{frame}



%%%%%%%%%%%%%%
\begin{frame}
\frametitle{Download data from Google Earth Engine}

\begin{enumerate}
    \item Wrote \href{\ghLtwo/python_demo/Prompt-data-scraping-Alpha-Earth}{Prompt-data-scraping-Alpha-Earth.md} pasted into ``plan mode''
    %\begin{itemize}
    %    \item I had to give permission on a few steps, e.g. searching URLs.
    %\end{itemize}
    \item Claude generated the plan \href{\ghLtwo/python_demo/alpha-earth-data-witty-planet.md}{alpha-earth-data-witty-planet.md} ($\sim$10 mins)
    \item Executed the plan to create \href{\ghLtwo/python_demo/download_satellite_embedding_v0.py}{download\_satellite\_embedding\_v0.py}
    \item Ran the code without edits (using Jupyter window).  issues: 
    \begin{itemize}
        \item Needed to add my GEE project code (ee-tannerregan-gwu)
        \item More substantial errors appeared, I iterated with Claude to make \href{\ghLtwo/python_demo/download_satellite_embedding_v1.py}{*\_v1.py}
    \end{itemize}
    \item Ran the updated script successfully
    \item Asked Claude to revise the code now that we had iterated
    \begin{itemize}
        \item Simplify code lines from the debugging/iterating process and make sure well commented \href{\ghLtwo/python_demo/download_satellite_embedding.py}{download\_satellite\_embedding.py}
    \end{itemize}
\end{enumerate}
\pause Check out VS Code Diff again!
\end{frame}



%%%%%%%%%%%%%%
\begin{frame}
\frametitle{Build a sample frame with visualization}
\begin{itemize}
    \item Wrote out three prompts (took me longer to get started with this exercise)
    \begin{itemize}
        \item \href{\ghLtwo/python_demo/Prompt-prepare-sample-frame.md}{Prompt-prepare-sample-frame.md} (writes draft)
        \begin{itemize}
            \item wrote plan: \href{\ghLtwo/python_demo/make-a-plan-based-zippy-sprout.md}{make-a-plan-based-zippy-sprout.md} 
        \end{itemize}
        \item \href{\ghLtwo/python_demo/Prompt-round2-prepare-sample-frame.md}{Prompt-round2-prepare-sample-frame.md} (adjustments)
        \begin{itemize}
            \item lost track of this plan (sorry)
        \end{itemize}
        \item \href{\ghLtwo/python_demo/Prompt-round3-prepare-sample-frame.md}{Prompt-round3-prepare-sample-frame.md} (added one new step)
        \begin{itemize}
            \item wrote plan: \href{\ghLtwo/python_demo/create-a-plan-that-witty-cosmos.md}{create-a-plan-that-witty-cosmos.md} 
        \end{itemize}
    \end{itemize}
    \item The final python script is \href{\ghLtwo/python_demo/prepare_sample_frame.py}{prepare\_sample\_frame.py}
\end{itemize}

\end{frame}

%%%%%%%%%%%%%%
\begin{frame}
\frametitle{Test an event study with satellite embeds}
\begin{itemize}
    \item Wrote a prompt to run a couple example designs
    \begin{itemize}
        \item \href{\ghLtwo/python_demo/Prompt-event-study.md}{Prompt-event-study.md} was my initial prompt
        \item This time I did many iterations and didn't keep track (sorry)
    \end{itemize}
    \item The final python script was \href{\ghLtwo/python_demo/event_study.py}{event\_study.py}
    \begin{itemize}
        \item Claude also make the  \href{\ghLtwo/python_demo/4_slides_event_study.tex}{4\_slides\_event\_study.tex} that built the slides at the top of this section
    \end{itemize}
\end{itemize}
\end{frame}


%%%%%%%%%%%%%%
\begin{frame}
\frametitle{Implications for RAs and research students}

\begin{itemize}
    \item My willingness to pay for PhD and MS student RAs is down
    \item Conditional on working with an RA, task allocation will change:
    \begin{itemize}
        \item $\downarrow$ More menial tasks like ugrad RA {\color{gray}(data collection, subjective classification, etc.)}
        \item $\uparrow$  More high level conceptual tasks {\color{gray}(design survey questionnaires, design data cleaning / empirical strategy)}
    \end{itemize} 
    \item This has implications for students 
    \begin{itemize}
        \item They still need to know the fundamentals of coding to make CC productive, but training opportunities will become more scarce
        \item Otherwise, RAships that focus more on high level conceptual tasks will probably help make them better researchers. But IMO requires higher RA ability and trust
    \end{itemize}
\end{itemize}

\end{frame}



\section{\textcolor{violet}{Wrapping Up \vspace{0.3in}}}


%%%%%%%%%%%%%%
\begin{frame}
\frametitle{Key Takeaways from Sessions 1 \& 2}

\begin{enumerate}
    \item AI handles the coding mechanics --- your job is to verify the output and drive the analysis
    \smallskip
    \item Use Plan mode for any non-trivial task; it gives you an audit trail and catches mistakes early
    \smallskip
    \item Git + Claude Code = version-controlled, reproducible research without learning Git commands
    \smallskip
    \item CLAUDE.md files carry your preferences across sessions so you do not repeat yourself
    \smallskip
    \item Replication packages are now a hours-long task, not a weeks-long one
    \smallskip
    \item The critical skill is \textit{verification}: check outputs, not just that the code runs
\end{enumerate}

\end{frame}


%%%%%%%%%%%%%%
\begin{frame}
\frametitle{Things to Try Before the Next Session}

\begin{itemize}
    \item Set up VS Code with Claude Code, Git, and Python (use the setup guide)
    \item Ask Claude to write a small data task in your preferred language --- verify the output
    \item Initialize a Git repo for one of your current projects
    \item Write a CLAUDE.md file for that project
    \vspace{.05in}
    \item Coming up in Session 3: Finding and Collecting Data
        \begin{itemize}
            \item[$\circ$] Identifying datasets, scraping, APIs, parsing PDFs at scale
        \end{itemize}
\end{itemize}

\end{frame}



%%%%%%%%%%%%%%
\begin{frame}
\frametitle{Questions?}

\begin{itemize}
    \item Questions on anything from Sessions 1 or 2?
    \item What tasks in your own research would you most like to try this on?
\end{itemize}

\end{frame}


\end{document}
